\begin{lemma}[The total mass of the subordination weight]
The function
\[
  s \longmapsto \frac{e^{-s}}{\sqrt s}
\]
is Lebesgue integrable on $(0,\infty)$, and
\[
  \int_0^\infty \frac{e^{-s}}{\sqrt s}\,ds = \sqrt\pi .
\]

This is the scalar normalization constant of the subordination formula
$L^{-1/2}\Pi_0 = \pi^{-1/2}\int_0^\infty s^{-1/2}P_s\Pi_0\,ds$ used throughout
the tablet's kernel layer: it is the total mass of the weight $s^{-1/2}e^{-s}$
against which every time integral in that layer is dominated, and it is the
reason the prefactor in the subordination formula is $\pi^{-1/2}$ and not some
other constant.  Equivalently it is the classical evaluation
$\Gamma(1/2) = \sqrt\pi$.  No Gaussian measure on $\mathbb{R}^n$ occurs here:
the statement is a one-dimensional Lebesgue integral of an explicit scalar
function, and it is recorded as its own node because several nodes of the kernel
layer need it as a majorant and as a normalization.
\end{lemma}
\begin{proof}
Write $g(s) = s^{-1/2}e^{-s}$ for $s > 0$.  The function $g$ is continuous on
$(0,\infty)$, hence Borel measurable there, and it is nonnegative.

\emph{Step 1: the substitution $s = v^2$.}
Let $\psi : (0,\infty) \to (0,\infty)$, $\psi(v) = v^2$.  It is a bijection of
$(0,\infty)$ onto $(0,\infty)$ with inverse $s \mapsto \sqrt s$, it is
continuously differentiable, and $\psi'(v) = 2v > 0$ on $(0,\infty)$.  By the
one-dimensional change-of-variables formula for a continuously differentiable
bijection with nonvanishing derivative, applied to the nonnegative Borel
function $g$,
\begin{equation*}
  \int_{(0,\infty)} g(s)\,ds
    = \int_{(0,\infty)} g(\psi(v))\,\psi'(v)\,dv , \tag{1}
\end{equation*}
both sides being elements of $[0,+\infty]$.

For $v > 0$ we have $\sqrt{\psi(v)} = \sqrt{v^2} = v$ because $v > 0$, so
\[
  g(\psi(v))\,\psi'(v)
    = \frac{e^{-v^2}}{v}\cdot 2v
    = 2\,e^{-v^2} .
\]
Hence $(1)$ reads
\begin{equation*}
  \int_{(0,\infty)} \frac{e^{-s}}{\sqrt s}\,ds
    = 2\int_{(0,\infty)} e^{-v^2}\,dv . \tag{2}
\end{equation*}

\emph{Step 2: evaluating the half-line Gaussian integral.}
The set $\{0\}$ is Lebesgue-null, so
\[
  \int_{\mathbb{R}} e^{-v^2}\,dv
    = \int_{(-\infty,0)} e^{-v^2}\,dv + \int_{(0,\infty)} e^{-v^2}\,dv ,
\]
the integrand being nonnegative and measurable, so that the additivity of the
integral over a partition into measurable pieces applies in $[0,+\infty]$.  By
the invariance of Lebesgue measure under the reflection $v \mapsto -v$ (shared
context), and since $e^{-(-v)^2} = e^{-v^2}$, the two half-line integrals are
equal.  By the scalar Gaussian integral (shared context),
$\int_{\mathbb{R}} e^{-v^2}\,dv = \sqrt\pi$.  Therefore
\[
  \int_{(0,\infty)} e^{-v^2}\,dv = \frac{\sqrt\pi}{2} .
\]

\emph{Step 3: conclusion.}
Substituting into $(2)$,
\[
  \int_{(0,\infty)} \frac{e^{-s}}{\sqrt s}\,ds = 2\cdot\frac{\sqrt\pi}{2} = \sqrt\pi .
\]
In particular this integral of the nonnegative measurable function $g$ is
finite, so $g$ is Lebesgue integrable on $(0,\infty)$, and its integral has the
stated value.
\end{proof}
